\section{Retroalimentación y automejora iterativa}

\subsection{Trace feedback y evaluación por pasos}
Chen presenta el trace feedback de AlphaCode, que usa lenguaje natural para reify la traza de ejecución del modelo, formando un stepwise verifier. Al hacer que el modelo genere su propio line-by-line debug, este puede identificar qué estados intermedios son promising, y usa una evaluación text-level que influye en el search posterior.

\begin{figure}[H]
\centering
\includegraphics[width=0.92\textwidth]{frame-trace-feedback.jpg}
\caption{Captura de ejemplo de trace feedback, que muestra la visualización de la line-by-line explanation.\protect\footnotemark}
\label{fig:trace-feedback}
\end{figure}
\footnotetext{Intervalo de tiempo en el video: 01:11:20--01:12:15, donde se explica el funcionamiento del trace feedback.}

\begin{importantbox}{El valor de más alto nivel del trace feedback}
1. Convierte la step evaluation en un natural language trace, haciendo que el propio modelo se convierta en verifier; 2. AlphaCode observa la context renovation a través del trace, lo que le permite detectar subtle bugs durante la etapa de code generation; 3. Un trace más rico puede usarse directamente como una signal de recompensa de RL, lo que ofrece una vía de mejora en inference-time para sistemas como AlphaProof y AlphaGeometry.
\end{importantbox}

\subsection{Prompt de autocorrección y riesgos}
Cuando no hay un oracle verifier, la autocorrección a menudo transforma un error en otro error distinto; Chen advierte de esto usando un ejemplo negativo en 01:13:23. Al diseñar un general-purpose feedback prompt, es necesario resaltar la consistency, incorporar el trace context y establecer guard rails del tipo «no alterar las partes ya confirmadas».

\begin{knowledgebox}{Direcciones para ajustar el prompt de autocorrección}
1. Resaltar la consistency: indicarle al modelo que «enumere dos posibles pasos nuevos», en lugar de una guía ambigua como «revisa la respuesta»; 2. Inyectar trace context: usar de nuevo la line-by-line explanation como prompt, para ayudar al modelo a juzgar la correctness con su propio reasoning; 3. Controlar los saltos: agregar la instrucción «conservar el razonamiento ya confirmado» para evitar que el modelo, por feedback drift, revierta las partes ya verificadas.
\end{knowledgebox}

\begin{warningbox}{Desventajas del self-correction sin oracle}
Cuando solo hay general feedback y no hay ground truth, el modelo, al no poder distinguir con claridad qué response es confiable, termina reescribiendo de forma conservadora una respuesta correcta o abandonando el reasoning de varios pasos, lo que provoca que la accuracy disminuya en vez de mejorar.
\end{warningbox}

\subsection{Resumen de la sección}
El trace feedback y el self-correction necesitan un reliable verifier y guard rails en el prompt; de lo contrario, el modelo retrocede después de varios rounds. Al diseñar el feedback hay que atender a la vez a la guía de consistency, al trace context y a los límites de interpretabilidad.
